\documentclass[12pt,a4paper]{article}
\usepackage[utf8]{inputenc}
\usepackage[T1]{fontenc}
\usepackage{amsmath,amssymb,amsthm}
\usepackage{geometry}
\usepackage[colorlinks=true,linkcolor=blue,citecolor=blue,urlcolor=blue]{hyperref}
\usepackage{listings}
\usepackage{xcolor}
\usepackage{booktabs}
\usepackage{enumitem}
\usepackage{mdframed}
\usepackage{mathtools}
\usepackage{microtype}
\usepackage{parskip}

\geometry{margin=2.5cm}

\lstset{
  basicstyle=\small\ttfamily,
  breaklines=true,
  frame=single,
  language=Mathematica,
  keywordstyle=\color{blue!70!black},
  commentstyle=\color{gray},
  stringstyle=\color{green!50!black},
  backgroundcolor=\color{gray!5},
  xleftmargin=4pt,
  xrightmargin=4pt,
  extendedchars=true,
  inputencoding=utf8,
  literate={->}{{\texttt{->}}}2
           {<|}{{\texttt{<|}}}2
           {|>}{{\texttt{|>}}}2
}

\newmdenv[
  backgroundcolor=blue!4,
  linecolor=blue!40,
  linewidth=1pt,
  innerleftmargin=10pt,
  innerrightmargin=10pt,
  innertopmargin=8pt,
  innerbottommargin=8pt
]{infobox}

\newmdenv[
  backgroundcolor=orange!8,
  linecolor=orange!60,
  linewidth=1pt,
  innerleftmargin=10pt,
  innerrightmargin=10pt,
  innertopmargin=8pt,
  innerbottommargin=8pt
]{warnbox}

\newmdenv[
  backgroundcolor=green!6,
  linecolor=green!50!black,
  linewidth=1pt,
  innerleftmargin=10pt,
  innerrightmargin=10pt,
  innertopmargin=8pt,
  innerbottommargin=8pt
]{notebox}

\title{\textbf{Towards a General Interface for SZ-DBC}\\[4pt]
\large Implementation Plan for Dimension-Agnostic, Automated\\
Detailed Balance Verification of Markov Chain Algorithms}
\author{S.~Whitby}
\date{June 2026}

\begin{document}
\maketitle
\tableofcontents
\newpage

% ============================================================
\section{Introduction and Motivation}

The current SZ-DBC checker verifies detailed balance for a single, fixed system
configuration: one specific lattice size, one specific particle multiset, one
specific spatial dimension (2D). These parameters are all hard-coded in the
algorithm file itself. This design conflates two concerns that should be
separate: the \emph{physics} of the algorithm under test, and the
\emph{infrastructure} needed to run the test.

This report describes a comprehensive plan to give SZ-DBC a general,
dimension-agnostic interface. The goal is threefold:

\begin{enumerate}
  \item An algorithm file should contain only the essential physics: the spatial
    dimension, the particle types, the transition rule, the energy function, and
    the coupling atoms. Everything else --- seed state construction, lattice
    size, particle count --- is generated automatically by the checker.

  \item The checker should sweep over all relevant system sizes and particle
    counts up to user-specified maxima, running the full BFS and DB
    verification at each point and aggregating the results into a single
    PASS/FAIL verdict.

  \item The framework should extend naturally to higher-dimensional systems,
    including \emph{phase-space} algorithms where each particle carries both a
    position and a momentum (or velocity), enabling verification of a richer
    class of physically important algorithms such as lattice kinetic Monte
    Carlo methods.
\end{enumerate}

Section~\ref{sec:architecture} describes the three-layer architecture.
Section~\ref{sec:interface} specifies the new algorithm file interface.
Section~\ref{sec:sweep} describes the automated seed-state sweep.
Section~\ref{sec:ddim} covers the $D$-dimensional generalisation of the core
library.  Section~\ref{sec:phasespace} addresses phase-space algorithms with
position and momentum, including the design of an archetypal example algorithm.
Section~\ref{sec:issues} catalogues foreseeable implementation issues.
Section~\ref{sec:effort} gives a difficulty and effort assessment.

% ============================================================
\section{Architecture}
\label{sec:architecture}

The reformed codebase has three clearly separated layers, each with a
well-defined responsibility.

\begin{center}
\begin{tabular}{llp{7cm}}
\toprule
Layer & File & Owns \\
\midrule
Algorithm & \texttt{my\_algo.wl} & physics: energy, transition rule, coupling atoms, spatial dimension \\
Orchestrator & \texttt{check.wls} & sweep strategy: which sizes and particle counts to test \\
Core library & \texttt{dbc\_core.wl} & generic machinery: BFS, orbit computation, DB check \\
\bottomrule
\end{tabular}
\end{center}

\medskip
Nothing from the middle row should appear in the top row, and nothing
algorithm-specific should be in the bottom two rows.  This is precisely what
the current design violates: every algorithm file sets \texttt{\$nGrid},
\texttt{\$particleTypes}, and \texttt{\$seedState}, and the core library
hard-codes 2D position vectors throughout.

The two documents that define the interface boundary are:
\begin{itemize}
  \item \texttt{algorithm\_requirements.md} --- the contract every algorithm
    file must satisfy.
  \item \texttt{template.wl} --- a minimal working algorithm file that a user
    copies and fills in.
\end{itemize}

% ============================================================
\section{Algorithm File Interface}
\label{sec:interface}

\subsection{Mandatory declarations}

\paragraph{\texttt{\$AlgDim = D}.}
The spatial dimension as a \textbf{hard-coded integer} --- the only
non-symbolic, non-overridable value in the algorithm file.  It cannot be made
symbolic because it determines the shape of all seed states, the structure of
the translation group, and the form of the $\tau$-BFS offset.  For a pure 2D
system \texttt{\$AlgDim = 2}; for a 2D position $+$ 2D momentum system, still
\texttt{\$AlgDim = 2} (the spatial dimension).

\paragraph{\texttt{\$particleTypeSet = \{$t_1, t_2, \ldots$\}}.}
The set of distinct particle types the algorithm recognises.  The checker will
generate all canonical type multisets from this set for each particle count
tested.

\paragraph{\texttt{Algorithm[state\_]}.}
The MCMC transition rule.  Must use \texttt{\$nGrid} as a \emph{free variable}
--- a symbol whose value is set by the checker before each run.  Must not
hard-code any specific lattice size internally.

\paragraph{\texttt{energy[state\_]}.}
The bare energy (no $\beta$ factor), with the same free-variable convention for
\texttt{\$nGrid}.

\paragraph{\texttt{DynamicSymParams[states\_]}.}
Returns the association of symbolic coupling atoms:
\begin{lstlisting}
<| "couplings" -> {couplingJ[1,2,1], couplingJ[1,2,2], ...},
   "extraSymParams" -> {} |>
\end{lstlisting}
Called once per run with the current state set.

\subsection{Optional declarations}

\paragraph{\texttt{\$symmetryGroup}.}
Defaults to \texttt{\{"translation"\}}.  The checker verifies any declared
symmetry before exploiting it.

\paragraph{\texttt{\$periodicDims}.}
Which coordinate components have periodic boundary conditions.  Defaults to
\texttt{Range[\$AlgDim]} (all spatial components).  For a phase-space system
with 4D coordinates $\{x, y, v_x, v_y\}$, set
\texttt{\$periodicDims = \{1, 2\}} so that only the spatial components receive
the $\tau$-offset and modular normalisation.  This is the mechanism by which
the framework generalises to phase-space algorithms without structural changes
to the BFS engine.

\paragraph{\texttt{\$velocitySet}.}
For phase-space algorithms: the finite set of allowed velocity vectors,
e.g.\ \texttt{\{\{-1,0\},\{1,0\},\{0,-1\},\{0,1\}\}}.  The checker uses this
to enumerate the full phase-space site set $\{1,\ldots,n\}^{\texttt{\$AlgDim}}
\times \texttt{\$velocitySet}$ when generating states.

\paragraph{\texttt{numBeta}.}
Inverse temperature for numerical mode only.

\subsection{What must NOT appear}

The algorithm file must \emph{not} set \texttt{\$nGrid} to a specific value,
define \texttt{\$seedState}, or define \texttt{\$particleTypes}.  These are
generated by the checker.

\begin{warnbox}
\textbf{Backward compatibility.} If the checker detects that the algorithm
file sets \texttt{\$nGrid}, it runs in \emph{single-size mode}: the declared
\texttt{\$nGrid}, \texttt{\$particleTypes}, and \texttt{\$seedState} are used
as-is, and the sweep is skipped.  All existing algorithm files in the
\texttt{examples/} directory work unchanged.
\end{warnbox}

\subsection{Minimal template}

\begin{lstlisting}
(* ==== template.wl -- minimal algorithm file ====
   Copy this file and fill in the three required sections. *)

$AlgDim         = 2;                   (* ONLY hard-coded value *)
$particleTypeSet = {1, 2};             (* distinct particle types *)
$symmetryGroup  = {"translation", "D4"};

(* ---- Energy ---- *)
energy[state_] := Module[{...},
  (* Use $nGrid as a free variable for lattice size.
     All spatial operations must use pairwise differences. *)
  ...]

(* ---- Algorithm ---- *)
Algorithm[state_List] := Module[{n = $nGrid, ...},
  (* One MCMC step. Call RandomReal[], RandomChoice[], etc. as needed.
     Do NOT hard-code n = 3 or any specific lattice size. *)
  ...]

(* ---- Symbolic parameters ---- *)
DynamicSymParams[states_List] :=
  Module[{types, couplings},
    types = Sort @ DeleteDuplicates @ Flatten[#[[2]] & /@ # & /@ states, 1];
    couplings = Flatten @ Table[
      couplingJ[a, b], {a, types}, {b, a, Max[types]}];
    <| "couplings" -> couplings, "extraSymParams" -> {} |>]
\end{lstlisting}

% ============================================================
\section{Automated Seed-State Sweep}
\label{sec:sweep}

\subsection{New command-line flags}

\begin{center}
\begin{tabular}{ll}
\toprule
Flag & Meaning \\
\midrule
\texttt{-maxN N} & Maximum lattice side length (default 3) \\
\texttt{-maxParticles K} & Maximum particles per seed state (default 4) \\
\texttt{-allTypes} & Test all canonical type multisets (default: one per
count) \\
\bottomrule
\end{tabular}
\end{center}

\subsection{Sweep structure}

For each lattice side length $n = 1, 2, \ldots, \texttt{maxN}$:
\begin{enumerate}
  \item Compute the number of available sites:
    $N_{\text{sites}} = n^{\texttt{\$AlgDim}} \times |\texttt{\$velocitySet}|$
    (the $|\texttt{\$velocitySet}|$ factor is 1 for pure-position systems).
  \item For each particle count $k = 0, 1, \ldots, \min(\texttt{maxParticles},
    N_{\text{sites}})$:
  \begin{enumerate}
    \item Generate one or more canonical type multisets of size $k$ from
      \texttt{\$particleTypeSet} (see below).
    \item For each such multiset $T$: set \texttt{\$nGrid} $= n$,
      \texttt{\$particleTypes} $= T$, generate the canonical seed state, and
      run the full Steps~1--8 pipeline.
    \item Record \texttt{PASS/FAIL} for the triple $(n, k, T)$.
  \end{enumerate}
\end{enumerate}
After all runs: print an aggregated table; exit with code 0 if all passed,
code 1 otherwise.

\subsection{Canonical type multisets}

Given particle count $k$ and type set $\{t_1, t_2, \ldots\}$, the default
canonical multiset cycles through types in order:
$T = (t_1, t_2, t_1, t_2, \ldots)$ truncated to length $k$.  For two types
$\{1,2\}$ and $k = 5$: $T = (1, 2, 1, 2, 1)$.

\begin{infobox}
\textbf{Why this default matters.}  With $k$ particles cycling through all
available types, every pair of distinct types appears in the system for even
$k$.  This ensures that all coupling atoms $\texttt{couplingJ}[a,b,d^2]$ can
appear in the transition probabilities, making the DB check sensitive to all
coupling parameters.  The \texttt{-allTypes} flag generates every distinct
multiset for thorough checking.
\end{infobox}

\subsection{Canonical seed state}

Given lattice side length $n$, dimension $D$, and type multiset $T = (t_1,
\ldots, t_k)$, the canonical seed state places particle $i$ at the $i$-th
position in lexicographic order on $\{1,\ldots,n\}^D$:
\begin{align*}
  \text{seedState} &= \bigl\{ \{(1,1,\ldots,1),\, t_1\},\;
                               \{(1,1,\ldots,2),\, t_2\},\; \ldots \bigr\}
\end{align*}
For a 2D lattice of side 3 with $T = (1, 2)$: seed $=\{((1,1),1),\;
((1,2),2)\}$.  This is deterministic, dimension-agnostic, and independent of
any physics.

For phase-space systems, the seed also specifies initial velocities: each
particle is assigned the first element of \texttt{\$velocitySet} (e.g.,
$(-1,-1)$), giving the lexicographically first phase-space state.

\subsection{Output format}

The checker prints a table after completing all runs:

\begin{lstlisting}
n=1, k=0, T={}         : tau PASS  D4 --   Erg --   DB PASS
n=1, k=1, T={1}        : tau PASS  D4 --   Erg PASS DB PASS
n=2, k=0, T={}         : tau PASS  D4 PASS Erg --   DB PASS
n=2, k=1, T={1}        : tau PASS  D4 PASS Erg PASS DB PASS
n=2, k=2, T={1,2}      : tau PASS  D4 PASS Erg PASS DB PASS
n=3, k=2, T={1,2}      : tau PASS  D4 PASS Erg PASS DB PASS
n=3, k=3, T={1,2,1}    : tau PASS  D4 PASS Erg PASS DB PASS
---
OVERALL: PASS (7/7 configurations)
\end{lstlisting}

% ============================================================
\section{\texorpdfstring{$D$}{D}-Dimensional Generalisation of the Core Library}
\label{sec:ddim}

\subsection{State format}

The current format \texttt{\{\{row, col\}, type\}} is replaced by
\texttt{\{\{$c_1, c_2, \ldots, c_{D_\text{total}}$\}, type\}} where
$D_\text{total} = \texttt{\$AlgDim} + \texttt{\$MomDim}$.  For a pure 2D
system $D_\text{total} = 2$; for a 2D spatial + 2D velocity system
$D_\text{total} = 4$.

This change is backwards-compatible: any function that currently indexes
\texttt{state[[i,1,1]]} (row) and \texttt{state[[i,1,2]]} (column) continues
to work for $D = 2$.  The internal functions must be generalised to loop over
the coordinate vector.

\subsection{State normalisation}

\begin{lstlisting}
(* Current (hardwired 2D) *)
$dbcNormState[state_, nGrid_] :=
  SortBy[Map[{{Mod[#[[1,1]]-1,nGrid]+1,
               Mod[#[[1,2]]-1,nGrid]+1}, #[[2]]} &, state], First]

(* Generalised: periodicDims specifies which components wrap *)
$dbcNormState[state_, nGrid_, periodicDims_] :=
  SortBy[Map[Function[p,
    {MapIndexed[
       If[MemberQ[periodicDims,#2[[1]]], Mod[#1-1,nGrid]+1, #1] &,
       p[[1]]],
     p[[2]]}], state], First]
\end{lstlisting}

The \texttt{periodicDims} parameter controls which coordinate components are
wrapped by the modular arithmetic.  For pure position systems it equals
\texttt{Range[AlgDim]}; for phase-space systems it is the subset
\texttt{\{1,\ldots,AlgDim\}}.

\subsection{\texorpdfstring{$\tau$}{Tau}-BFS in \texorpdfstring{$D$}{D} spatial dimensions}

The translational invariance check adds symbolic offsets $\tau_1, \ldots,
\tau_D$ to the spatial components of every position vector.  Generically, for
spatial dimension $D$, the checker introduces $D$ symbols \texttt{\{$\tau_1$,
\ldots, $\tau_D$\}} and augments only the periodic components:

\begin{lstlisting}
(* Generalised tau-augmentation *)
$dbcAddTau[state_, alg_dim_] :=
  Map[Function[p,
    {MapIndexed[
       If[#2[[1]] <= alg_dim, #1 + Symbol["tau"<>ToString[#2[[1]]]], #1] &,
       p[[1]]],
     p[[2]]}], state]
\end{lstlisting}

For a 2D system this recovers the current $\{\tau_r, \tau_c\}$ behaviour.  For
a phase-space system with coordinates $\{x, y, v_x, v_y\}$ and
\texttt{\$AlgDim = 2}, the augmented positions are $\{x + \tau_1, y + \tau_2,
v_x, v_y\}$ --- velocity components are not shifted, correctly capturing that
translational invariance is a property of position space alone.

\subsection{Translation group in \texorpdfstring{$D$}{D} dimensions}

The group of lattice translations $\mathbb{Z}_n^D$ is generated by all shifts
$\mathbf{d} = (d_1, \ldots, d_D)$ with $d_i \in \{0, 1, \ldots, n-1\}$.
Currently the group elements are enumerated by a 2D \texttt{Table} over
\texttt{\{dr, dc\}}; the generalisation is a $D$-fold Cartesian product:

\begin{lstlisting}
(* Generalised translation elements *)
translationElems[nGrid_, alg_dim_] :=
  Map[<| "shift" -> # |> &,
    Flatten[Outer[List, Sequence @@ Table[Range[0,nGrid-1], alg_dim]], alg_dim-1]]
\end{lstlisting}

The shift vectors act on the spatial components of position only; velocity
components are unchanged.

\subsection{Rotational symmetry}

For $D = 2$: the D4 group (8 elements) is already implemented.  No change.

For $D = 3$: the octahedral group $O_h$ has 48 elements.  These can be
enumerated as the set of all $3\times3$ signed permutation matrices with
determinant $+1$.  Adding Oh to the checker is a well-defined but non-trivial
extension; it is deferred to a later phase.

For phase-space systems ($D = 2$ spatial, $D = 2$ velocity): a D4 rotation by
$90^\circ$ clockwise maps $(x, y) \to (y, n+1-x)$ in position space and
simultaneously maps $(v_x, v_y) \to (v_y, -v_x)$ in velocity space.  This is
a joint transformation on the full coordinate vector.  The symmetry element
must encode both actions:

\begin{lstlisting}
(* D4 rotation for a 2D-position + 2D-velocity state *)
$dbcPhaseSpaceD4Pos[1, x_, y_, vx_, vy_, n_] :=
  {y, n+1-x, vy, -vx}   (* 90 degrees CW *)
\end{lstlisting}

For the initial implementation of phase-space systems, translation-only orbits
are sufficient.  The coupled D4 action can be added in a later phase.

\subsection{State enumeration}

\begin{lstlisting}
(* Current: 2D grid positions only *)
$dbcEnumerateNParticleStates[seedState_, nGrid_] :=
  Module[{types, allPos},
    types  = Sort[#[[2]] & /@ seedState];
    allPos = Flatten[Table[{r,c},{r,nGrid},{c,nGrid}],1];
    ...]

(* Generalised: D-dimensional grid, optional velocity set *)
$dbcEnumerateNParticleStates[types_, nGrid_, alg_dim_, velocity_set_:{}] :=
  Module[{posSites, sites, allAssignments},
    posSites = Flatten[
      Outer[List, Sequence @@ Table[Range[nGrid], alg_dim]],
      alg_dim - 1];
    sites = If[velocity_set === {},
      posSites,
      Flatten[Outer[Join, posSites, velocity_set, 1], 1]];
    (* Enumerate all N-particle assignments on `sites` *)
    ...]
\end{lstlisting}

For a 2D lattice of side $n$ with 4-direction velocities, this produces $4n^2$
sites.  For $n = 2$ and $N = 2$ particles: $C(16, 2) = 120$ two-particle
states (assuming particles are distinguishable by type).

% ============================================================
\section{Phase-Space Algorithms}
\label{sec:phasespace}

\subsection{Why phase-space matters}

Most existing MCMC algorithms in statistical physics evolve only the positions
of particles; velocities or momenta are integrated out analytically.  A richer
class of algorithms, including stochastic lattice gases and kinetic Monte Carlo
methods, maintains both positions and velocities as dynamical variables and
samples from the full phase-space Boltzmann distribution
\[
  \pi(\mathbf{r}, \mathbf{v}) \propto
    \exp\!\bigl(-\beta\, H(\mathbf{r}, \mathbf{v})\bigr), \qquad
  H = \sum_i \tfrac{1}{2}|\mathbf{v}_i|^2 + \sum_{i<j} V(|\mathbf{r}_i -
    \mathbf{r}_j|^2).
\]
Verifying detailed balance for these algorithms requires the checker to handle
states with a velocity component, which the current design does not support.

\subsection{State format for phase-space algorithms}

Each particle is described by a four-component coordinate vector $\{x, y, v_x,
v_y\}$ plus a type.  The full state is a sorted list of such pairs:
\begin{lstlisting}
{  {{1, 1, -1, 0}, 1},
   {{2, 3,  0, 1}, 2},
   ...  }
\end{lstlisting}
Sorting is by the full coordinate vector lexicographically: positions are
primary, velocities are secondary.  This generalises the current sort without
any conceptual change.

The algorithm file declares:
\begin{lstlisting}
$AlgDim       = 2;               (* x, y are the two spatial dims *)
$periodicDims = {1, 2};          (* only x, y are periodic *)
$velocitySet  = {{-1,0},{1,0},{0,-1},{0,1}};  (* allowed velocities *)
\end{lstlisting}

\subsection{The streaming--collision problem}
\label{sec:streamcol}

The most physically natural algorithm for phase-space systems is
streaming--collision dynamics:
\begin{enumerate}
  \item \textbf{Streaming}: each particle moves $\mathbf{r}_i \to \mathbf{r}_i
    + \mathbf{v}_i$ (mod $n$).
  \item \textbf{Collision}: at sites with multiple particles, redistribute
    velocities stochastically.
\end{enumerate}

However, constructing a streaming--collision algorithm that satisfies detailed
balance with the \emph{full} Hamiltonian $H = K + V$ is non-trivial.  The
streaming step is deterministic and maps the system from one configuration to
another with different potential energy $V(\mathbf{r} + \mathbf{v}) \neq
V(\mathbf{r})$.  For the collision acceptance probability to restore detailed
balance, it must account for \emph{both} the kinetic energy change $\Delta K$
and the potential energy change $\Delta V$ induced by the preceding streaming.
This couples the collision kernel to the global state in a way that makes the
acceptance criterion complex and the DB verification non-trivial to design.

Specifically, for the combined step to satisfy
\[
  T(\mathbf{s} \to \mathbf{s}')
  \exp(-\beta H(\mathbf{s})) =
  T(\mathbf{s}' \to \mathbf{s})
  \exp(-\beta H(\mathbf{s}')),
\]
the collision acceptance ratio must equal
$\exp(-\beta(H(\mathbf{s}') - H(\mathbf{s}))) = \exp(-\beta(\Delta K +
\Delta V))$, where $\Delta V$ depends on the \emph{post-streaming} particle
positions.  A collision kernel that uses only $\Delta K$ (the standard choice)
satisfies detailed balance with respect to $K$ alone --- which is correct only
when $V \equiv 0$ (non-interacting particles).

This is a genuine theoretical constraint, not a limitation of the checker.  Any
correct streaming--collision MCMC algorithm for interacting phase-space
particles must either: (a) use a collision acceptance criterion that references
the global potential energy, or (b) apply a separate global Metropolis
correction step after streaming.  Implementing such an algorithm cleanly is
deferred to a future extension.

\subsection{Archetypal phase-space algorithm: position--velocity Metropolis}
\label{sec:pvmetropolis}

The clean, pedagogically correct first example is a \emph{position--velocity
Metropolis} algorithm: two independent sub-moves, chosen at random, each of
which satisfies detailed balance with respect to the full Hamiltonian $H = K +
V$.

\subsubsection*{Algorithm description}

At each step:
\begin{enumerate}
  \item Choose a particle $p$ uniformly at random.
  \item Draw $u \sim \mathrm{Uniform}[0,1]$.
  \item \textbf{If $u < \tfrac{1}{2}$ (position move):}
    \begin{enumerate}
      \item Choose a displacement $\mathbf{d}$ uniformly from the set of
        nearest-neighbour directions.
      \item Proposed new position: $\mathbf{r}' = \mathbf{r}_p + \mathbf{d}$.
      \item \emph{Hard-core rejection}: if any other particle already occupies
        $\mathbf{r}'$ (regardless of its velocity), return the current state.
      \item Compute $\Delta V = V_\text{new} - V_\text{old}$ (potential energy
        change; velocity of $p$ is unchanged).
      \item Accept with $\min(1, e^{-\beta \Delta V})$; else stay.
    \end{enumerate}
  \item \textbf{If $u \geq \tfrac{1}{2}$ (velocity move):}
    \begin{enumerate}
      \item Choose a new velocity $\mathbf{v}'$ uniformly from
        \texttt{\$velocitySet}.
      \item Compute $\Delta K = \tfrac{1}{2}(|\mathbf{v}'|^2 - |\mathbf{v}_p|^2)$.
      \item Accept with $\min(1, e^{-\beta \Delta K})$; else stay.
    \end{enumerate}
\end{enumerate}

\subsubsection*{Why this satisfies detailed balance}

The position move satisfies DB with respect to $H$:
\[
  T_{\text{pos}}(\mathbf{s} \to \mathbf{s}')\, e^{-\beta H(\mathbf{s})} =
  T_{\text{pos}}(\mathbf{s}' \to \mathbf{s})\, e^{-\beta H(\mathbf{s}')}.
\]
This holds because (i) the proposal is symmetric (same neighbour set forward
and reverse), (ii) hard-core rejection is symmetric, and (iii) the Metropolis
acceptance satisfies $\min(1, e^{-\beta \Delta V}) \cdot e^{-\beta K} =
\min(1, e^{+\beta \Delta V}) \cdot e^{-\beta K} \cdot e^{-\beta \Delta V}$
exactly.  Kinetic energy cancels because $\Delta K = 0$ for a position move.

The velocity move satisfies DB with respect to $H$:
\[
  T_{\text{vel}}(\mathbf{s} \to \mathbf{s}')\, e^{-\beta H(\mathbf{s})} =
  T_{\text{vel}}(\mathbf{s}' \to \mathbf{s})\, e^{-\beta H(\mathbf{s}')}.
\]
This holds because the velocity proposal is symmetric over
\texttt{\$velocitySet}, potential energy cancels ($\Delta V = 0$ for a velocity
move), and Metropolis on $\Delta K$ satisfies the condition exactly.

A 50/50 random mixture of two Markov kernels that both satisfy DB with respect
to the same stationary distribution $\pi$ is itself DB-satisfying with respect
to $\pi$.  Therefore the full algorithm satisfies detailed balance with respect
to $\pi \propto \exp(-\beta H)$.

\subsubsection*{Algorithm file sketch}

\begin{lstlisting}
(* ==== phase_space_metropolis.wl ====
   2D spatial + 2D velocity Metropolis.
   State: {{{x, y, vx, vy}, type}, ...}
   H = sum_i (vx_i^2 + vy_i^2)/2 + sum_{i<j} couplingJ[ta,tb,|ri-rj|^2] *)

$AlgDim        = 2;
$periodicDims  = {1, 2};
$velocitySet   = {{-1,0},{1,0},{0,-1},{0,1}};
$particleTypeSet = {1, 2};
$maxD2         = 2;

$posDisps = {{-1,0},{1,0},{0,-1},{0,1}};

(* Potential energy: pairwise, depends on spatial positions only *)
$potE[state_, n_] := Module[{total = 0},
  Do[
    With[{d2 = $pairD2Spatial[state[[i,1]], state[[j,1]], n]},
      If[0 < d2 <= $maxD2,
        total += couplingJ[state[[i,2]], state[[j,2]], d2]]],
    {i, Length[state]}, {j, i+1, Length[state]}]; total]

$pairD2Spatial[p1_, p2_, n_] :=
  With[{dr = p1[[1]]-p2[[1]], dc = p1[[2]]-p2[[2]]},
    With[{dra = Mod[Abs[dr],n], dca = Mod[Abs[dc],n]},
      Min[dra,n-dra]^2 + Min[dca,n-dca]^2]]

(* Kinetic energy: sum of |v|^2 / 2 per particle *)
$kinE[state_] := Total[(#[[1,3]]^2 + #[[1,4]]^2)/2 & /@ state]

energy[state_] := $kinE[state] + $potE[state, $nGrid]

Algorithm[state_List] :=
  Module[{n = $nGrid, p, x, y, vx, vy, tp, rest, u},
    If[state === {}, Return[state]];
    p  = RandomChoice[state];
    x  = p[[1,1]]; y  = p[[1,2]];
    vx = p[[1,3]]; vy = p[[1,4]];
    tp = p[[2]];
    rest = DeleteCases[state, p, 1];
    u = RandomReal[];

    If[u < 1/2,
      (* --- Position move --- *)
      With[{d = RandomChoice[$posDisps],
            newXY = {x,y} + RandomChoice[$posDisps]},
        (* hard-core: no other particle at newXY regardless of velocity *)
        If[AnyTrue[rest, Function[q,
              Mod[q[[1,1]]-newXY[[1]],n]===0 &&
              Mod[q[[1,2]]-newXY[[2]],n]===0]],
           Return[state]];
        With[{ns = SortBy[Append[rest,
                    {{newXY[[1]],newXY[[2]],vx,vy},tp}], First]},
          With[{dV = $potE[ns,n] - $potE[state,n]},
            If[RandomReal[] < Piecewise[{{1,dV<=0}},Exp[-\[Beta]*dV]],
               ns, state]]]],

      (* --- Velocity move --- *)
      With[{nv = RandomChoice[$velocitySet]},
        With[{dK = (nv[[1]]^2+nv[[2]]^2 - vx^2-vy^2)/2},
          With[{ns = SortBy[Append[rest,{{x,y,nv[[1]],nv[[2]]},tp}],First]},
            If[RandomReal[] < Piecewise[{{1,dK<=0}},Exp[-\[Beta]*dK]],
               ns, state]]]]]]

DynamicSymParams[states_List] :=
  Module[{types, n, d2Vals, atoms},
    types  = Sort @ DeleteDuplicates @ Flatten[#[[2]] & /@ # & /@ states, 1];
    n      = $nGrid;
    d2Vals = Sort @ DeleteDuplicates @ Select[
      Flatten @ Table[Min[dr,n-dr]^2+Min[dc,n-dc]^2,
        {dr,0,Floor[n/2]},{dc,0,Floor[n/2]}],
      0 < # <= $maxD2 &];
    atoms = Flatten @ Table[
      If[a<=b, Table[couplingJ[a,b,d2],{d2,d2Vals}], Nothing],
      {a,types},{b,types}];
    <| "couplings" -> atoms, "extraSymParams" -> {} |>]
\end{lstlisting}

\subsubsection*{What the checker will verify}

When run with \texttt{-maxN 2 -maxParticles 2}, the checker will test:
\begin{itemize}
  \item $n=1$, $k=0$: trivially empty state.
  \item $n=1$, $k=1$: one particle on a single site; only velocity moves are
    possible.  The checker verifies DB for all velocity pairs $\{v, v'\}$ and
    all coupling values.
  \item $n=2$, $k=1$: one particle on a $2\times2$ torus.  Position moves and
    velocity moves both active.
  \item $n=2$, $k=2$, $T=\{1,2\}$: two particles on a $2\times2$ torus with 4
    velocities each.  State space has $C(16, 2) = 120$ states.  The DB check
    exercises all coupling atoms $\texttt{couplingJ}[1, 2, d^2]$ and verifies
    that the kinetic acceptance is correctly applied.
\end{itemize}

A broken variant (e.g., using $\beta/2$ in the velocity acceptance criterion)
will produce DB violations at the $k=1$ level where only velocity moves
operate, making the failure immediately legible.

% ============================================================
\section{Foreseeable Issues}
\label{sec:issues}

\subsection{Lazy evaluation of grid-dependent globals}
Several existing algorithm files compute geometry helpers at load time using
the then-current value of \texttt{\$nGrid}:
\begin{lstlisting}
$nStep       = Max[1, Round[physLen]];   (* vmmc_2d.wl *)
$smDisps     = DeleteCases[Flatten[...], {0,0}];  (* single_metropolis.wl *)
\end{lstlisting}
These are evaluated once when the file is \texttt{Get[]}'d and then cached as
global constants.  Under the sweep design, \texttt{\$nGrid} changes between
runs, so any helper that depends on \texttt{\$nGrid} and is computed at load
time will silently use a stale value.  The fix is to compute such helpers
lazily inside \texttt{Algorithm[]} or \texttt{energy[]} using the local
\texttt{n = \$nGrid} binding:
\begin{lstlisting}
Algorithm[state_List] := Module[{n = $nGrid,
    disps = DeleteCases[Flatten[Table[{dx,dy},{dx,-1,1},{dy,-1,1}],1],{0,0}],
    ...}, ...]
\end{lstlisting}
For the existing examples, inspection shows that the precomputed globals are
all independent of \texttt{\$nGrid} (e.g., \texttt{\$smDisps} is the fixed
list of 8 unit displacements).  The sweep design is therefore backward
compatible for these files without modification.

\subsection{EBE cost across many system sizes}
The EBE Phase~2 (hyperplane arrangement BFS via
\texttt{FindInstance}) runs once per $(n, T)$ combination.  A sweep of
\texttt{-maxN 3 -maxParticles 4} with two types produces roughly 20
independent check runs.  Phase~2 is fast for small systems but the aggregate
cost grows.  The existing chamber-reuse mechanism (D4 EBE chambers reused for
DB EBE) already eliminates duplicate Phase~2 work within a single $(n, T)$
run.  Cross-size reuse is not straightforward since the coupling atom set
changes with particle types.  The recommended mitigation is to fail fast: if
any $(n, T)$ run finds a DB violation, the sweep stops at that point.

\subsection{Phase-space state space growth}
Including velocity degrees of freedom multiplies the state space by
$|\texttt{\$velocitySet}|$ per particle.  For a $2 \times 2$ lattice with
4-direction velocities and 2 particles the state space is approximately
$120$, compared to $12$ for the pure-position case.  For $n = 3$ with 8
velocities and 3 particles the state space reaches several thousand.  For
phase-space systems the default \texttt{-maxN} should be 2, and the template
file should carry a comment recommending this.

\subsection{\texorpdfstring{$\tau$}{Tau}-BFS with four-dimensional coordinates}
For a state with coordinate vector $\{x, y, v_x, v_y\}$, the
\texttt{\$dbcAddTau} function must add $\tau_1$ to $x$ and $\tau_2$ to $y$
only.  If the \texttt{\$periodicDims} parameter is not passed through correctly
to all three $\tau$-related functions (\texttt{\$dbcAddTau},
\texttt{\$dbcNormTauState}, and the norm function passed to
\texttt{\$dbcBuildStateLeaves}), the $\tau$ symbols will either leak into the
velocity components (causing a spurious $\tau$-FAIL) or fail to cancel in the
spatial components (also a spurious FAIL).  Careful parameter threading is
required.

\subsection{Hard-core constraint for phase-space systems}
There are two natural hard-core conventions: (a) no two particles at the same
\emph{position}, regardless of velocity, and (b) no two particles at the same
\emph{phase-space point} $\{x, y, v_x, v_y\}$.  Convention (a) is physically
appropriate for particles with a finite size; convention (b) models a fermionic
exclusion principle.  These lead to different state enumerations and different
transition probabilities.  The algorithm file must make this choice explicit,
and the checker's state enumeration must respect it.  For the canonical example
(\S\ref{sec:pvmetropolis}) convention (a) is used because the position move
checks for occupancy based on position only.

\subsection{Coupling atoms with kinetic energy}
For a phase-space system the energy is $H = K + V$.  The kinetic part $K =
\sum_i |\mathbf{v}_i|^2/2$ uses the \emph{velocity values as data} from the
state, not as symbolic parameters.  The coupling atoms returned by
\texttt{DynamicSymParams} are therefore only the potential couplings
$\{\texttt{couplingJ}[a, b, d^2]\}$; the factor $1/2$ in $K$ is a fixed
rational.  This means \texttt{\$dbcIsExpZero} will encounter terms with both
rational exponents (from kinetic contributions) and symbolic exponents (from
potential contributions).  Fortunately \texttt{\$dbcIsExpZero} already handles
this: it groups terms by exponent class, and rational exponents form their own
class.  No changes to the DB check are needed.

% ============================================================
\section{Effort and Difficulty Assessment}
\label{sec:effort}

\begin{center}
\begin{tabular}{p{6.5cm}lp{4.5cm}}
\toprule
Task & Difficulty & Notes \\
\midrule
\texttt{algorithm\_requirements.md} + \texttt{template.wl} & Low &
  Documentation only \\
Backward-compat detection in \texttt{check.wls} & Low &
  Detect \texttt{\$nGrid} assignment \\
Multi-size sweep loop in \texttt{check.wls} & Medium &
  Nested loop; output table \\
Canonical type-multiset and seed-state generation & Medium &
  Deterministic, dimension-agnostic \\
$D$-dim state normalisation in \texttt{dbc\_core.wl} & Medium &
  Pervasive; thread \texttt{periodicDims} \\
$D$-dim translation group elements & Low-Medium &
  $D$-fold Cartesian product \\
$D$-dim orbit computation & Medium &
  Only positional indexing changes \\
$\tau$-BFS with \texttt{periodicDims} & Medium &
  Thread param into three functions \\
Phase-space state enumeration & Medium &
  Extend to velocity-set sites \\
\texttt{phase\_space\_metropolis.wl} example & Medium &
  Position and velocity moves clean \\
Coupled D4 on position+velocity & Hard &
  Defer to later phase \\
Streaming--collision MCMC & Hard &
  Non-trivial acceptance; defer \\
Adapt existing examples & Low &
  Remove 3 lines per file \\
Full regression testing & Medium &
  Verify sweep matches current output \\
\midrule
\textbf{Estimate (staged)} & & Core: 1--2 wks; phase-space example: 3--5 days \\
\bottomrule
\end{tabular}
\end{center}

\medskip
The recommended implementation order is:
\begin{enumerate}
  \item Write \texttt{algorithm\_requirements.md} and \texttt{template.wl}
    (establishes the interface contract before any code changes).
  \item Generalise \texttt{dbc\_core.wl} to $D$ dimensions with
    \texttt{periodicDims} (enables all subsequent work).
  \item Add the multi-size sweep to \texttt{check.wls} with backward-compat
    mode.
  \item Adapt existing examples and run regression tests.
  \item Implement \texttt{phase\_space\_metropolis.wl} and test on $n=2$.
  \item Add coupled D4 action and streaming--collision example (later phases).
\end{enumerate}

% ============================================================
\section{Summary}

The principal architectural change is a clean separation of \emph{physics}
(algorithm file) from \emph{infrastructure} (checker and core library).  An
algorithm file in the new design declares only three physical items: the spatial
dimension \texttt{\$AlgDim}, the particle type set \texttt{\$particleTypeSet},
and the three functions \texttt{Algorithm}, \texttt{energy}, and
\texttt{DynamicSymParams}.  Everything else --- lattice size, seed state,
particle count sweep --- is generated and controlled by the checker.

The $D$-dimensional generalisation of the core library is largely mechanical:
every function that currently indexes a 2-component position vector must be
updated to loop over the \texttt{periodicDims} parameter.  The BFS engine and
the DB checking machinery require no changes.

Phase-space algorithms are accommodated naturally within this framework by
treating velocity as additional (non-periodic) coordinate components.  The
archetypal example --- a position+velocity Metropolis algorithm with separate
sub-moves --- is theoretically clean, straightforward to implement, and a
valuable demonstration that the checker correctly handles a Hamiltonian with
both kinetic and potential terms.  The more complex streaming+collision case is
deferred because constructing a DB-correct streaming+collision algorithm for
\emph{interacting} particles requires an acceptance criterion that accounts for
potential energy changes due to streaming, making it a non-trivial algorithm
design problem in its own right.

\end{document}
